% GENERATED FILE. Edit canonical lessons, then run npm run generate:books.
% canonical-source-hash: fnv1a64:7d3bfdb98c82a00e
% canonical-lessons: SA-C248-kosah, SA-C248-nidhih, SA-C248-rupyakam, SA-C248-panah, SA-C248-sulkam

\chapter{Sheath, Treasure, Rupee, Stake, Fee}
\label{ch:sa-sheath-treasure-rupee-stake-fee}
\hlchaptermodality{248}

\begin{chapteropening}
\textbf{By the end of this chapter:} \emph{I can say a sheath, a treasure, a rupee, a stake and a fee.}

Complete the last lesson of chapter 248: I can say a sheath, a treasure, a rupee, a stake and a fee.
\end{chapteropening}

\section[\texorpdfstring{\sk{कोशः} (kośaḥ)}{कोशः (kośaḥ)}]{\sk{कोशः} (kośaḥ) — a sheath, a treasury}

\label{lesson:SA-C248-kosah}



\begin{quote}
Before the new one: say the Sanskrit for a flag, then the Sanskrit for a drum.
\end{quote}

\subsection*{You'll want to know: \sk{कोशः}}
\textbf{\sk{कोशः}} — \emph{kośaḥ} — \textquotedblleft{}a sheath, a treasury\textquotedblright{}.

\begin{grammarlens}[title={What you've built}]
One more word for this chapter.
\end{grammarlens}

\subsection*{Your turn}
\begin{itemize}
\raggedright
  \item \emph{Say it:} \emph{kośaḥ}
  \item \emph{Say it:} \emph{kośaḥ}, once more
  \item \emph{Say it:} say \emph{kośaḥ}
  \item \emph{Recall:} say the Sanskrit for a flag, then the Sanskrit for a drum, then say \emph{kośaḥ} again
\end{itemize}

\begin{tcolorbox}[breakable,colback=teal!4,colframe=teal!35!black,title={Before you move on}]
What does \sk{कोशः} mean? (\textquotedblleft{}A sheath, a treasury\textquotedblright{}.) Say it once more.
\end{tcolorbox}
\section[\texorpdfstring{\sk{निधिः} (nidhiḥ)}{निधिः (nidhiḥ)}]{\sk{निधिः} (nidhiḥ) — a treasure}

\label{lesson:SA-C248-nidhih}



\begin{quote}
Before the new one: say the Sanskrit for a drum, then the Sanskrit for a sheath.
\end{quote}

\subsection*{You'll want to know: \sk{निधिः}}
\textbf{\sk{निधिः}} — \emph{nidhiḥ} — \textquotedblleft{}a treasure\textquotedblright{}.

\begin{grammarlens}[title={What you've built}]
One more word for this chapter.
\end{grammarlens}

\subsection*{Your turn}
\begin{itemize}
\raggedright
  \item \emph{Say it:} \emph{nidhiḥ}
  \item \emph{Say it:} \emph{nidhiḥ}, once more
  \item \emph{Say it:} say \emph{nidhiḥ}
  \item \emph{Recall:} say the Sanskrit for a drum, then the Sanskrit for a sheath, then say \emph{nidhiḥ} again
\end{itemize}

\begin{tcolorbox}[breakable,colback=teal!4,colframe=teal!35!black,title={Before you move on}]
What does \sk{निधिः} mean? (\textquotedblleft{}A treasure\textquotedblright{}.) Say it once more.
\end{tcolorbox}
\section[\texorpdfstring{\sk{रूप्यकम्} (rūpyakam)}{रूप्यकम् (rūpyakam)}]{\sk{रूप्यकम्} (rūpyakam) — a rupee}

\label{lesson:SA-C248-rupyakam}



\begin{quote}
Before the new one: say the Sanskrit for a sheath, then the Sanskrit for a treasure.
\end{quote}

\subsection*{You'll want to know: \sk{रूप्यकम्}}
\textbf{\sk{रूप्यकम्}} — \emph{rūpyakam} — \textquotedblleft{}a rupee\textquotedblright{}.

\begin{grammarlens}[title={What you've built}]
One more word for this chapter.
\end{grammarlens}

\subsection*{Your turn}
\begin{itemize}
\raggedright
  \item \emph{Say it:} \emph{rūpyakam}
  \item \emph{Say it:} \emph{rūpyakam}, once more
  \item \emph{Say it:} say \emph{rūpyakam}
  \item \emph{Recall:} say the Sanskrit for a sheath, then the Sanskrit for a treasure, then say \emph{rūpyakam} again
\end{itemize}

\begin{tcolorbox}[breakable,colback=teal!4,colframe=teal!35!black,title={Before you move on}]
What does \sk{रूप्यकम्} mean? (\textquotedblleft{}A rupee\textquotedblright{}.) Say it once more.
\end{tcolorbox}
\section[\texorpdfstring{\sk{पणः} (paṇaḥ)}{पणः (paṇaḥ)}]{\sk{पणः} (paṇaḥ) — a stake, a wager, a coin}

\label{lesson:SA-C248-panah}



\begin{quote}
Before the new one: say the Sanskrit for a treasure, then the Sanskrit for a rupee.
\end{quote}

\subsection*{You'll want to know: \sk{पणः}}
\textbf{\sk{पणः}} — \emph{paṇaḥ} — \textquotedblleft{}a stake, a wager, a coin\textquotedblright{}.

\begin{grammarlens}[title={What you've built}]
One more word for this chapter.
\end{grammarlens}

\subsection*{Your turn}
\begin{itemize}
\raggedright
  \item \emph{Say it:} \emph{paṇaḥ}
  \item \emph{Say it:} \emph{paṇaḥ}, once more
  \item \emph{Say it:} say \emph{paṇaḥ}
  \item \emph{Recall:} say the Sanskrit for a treasure, then the Sanskrit for a rupee, then say \emph{paṇaḥ} again
\end{itemize}

\begin{tcolorbox}[breakable,colback=teal!4,colframe=teal!35!black,title={Before you move on}]
What does \sk{पणः} mean? (\textquotedblleft{}A stake, a wager, a coin\textquotedblright{}.) Say it once more.
\end{tcolorbox}
\section[\texorpdfstring{\sk{शुल्कम्} (śulkam)}{शुल्कम् (śulkam)}]{\sk{शुल्कम्} (śulkam) — a fee, a tax}

\label{lesson:SA-C248-sulkam}



\begin{quote}
Before the new one: say the Sanskrit for a rupee, then the Sanskrit for a stake.
\end{quote}

\subsection*{You'll want to know: \sk{शुल्कम्}}
\textbf{\sk{शुल्कम्}} — \emph{śulkam} — \textquotedblleft{}a fee, a tax\textquotedblright{}.

\begin{grammarlens}[title={What you've built}]
One more word for this chapter.
\end{grammarlens}

\subsection*{Your turn}
\begin{itemize}
\raggedright
  \item \emph{Say it:} \emph{śulkam}
  \item \emph{Say it:} \emph{śulkam}, once more
  \item \emph{Say it:} say \emph{śulkam}
  \item \emph{Recall:} say the Sanskrit for a rupee, then the Sanskrit for a stake, then say \emph{śulkam} again
\end{itemize}

\begin{tcolorbox}[breakable,colback=teal!4,colframe=teal!35!black,title={Before you move on}]
What does \sk{शुल्कम्} mean? (\textquotedblleft{}A fee, a tax\textquotedblright{}.) Say it once more.
\end{tcolorbox}
